\section{Monte Carlo Basic}
    \subsection{Main Idea}
        The \textbf{Monte Carlo basic} is the simplest Monte Carlo method algorithm.
        Suppose that we run the agent starting from state-action pair $(s,a)$, following policy $\pi$, and let $g_\pi^{(i)}(s,a)$
        be the return of the episode. Then, we can approximate the action value $q_\pi(s,a)$ as following:
        \begin{equation}
            q_\pi(s,a) = \E[G_t|S_t=s,A_t=a] \approx \frac{1}{n} \sum_{i=1}^n g_\pi^{(i)}(s,a).
        \end{equation}
        If the number $n$ is sufficiently large, the approximation will be sufficiently accurate. This can be guaranteed due to the law of large numbers.
    \subsection{Pseudocode}
        \Algorithm{mcBasic}
        {Monte Carlo Basic}
        {
            \inputitem{$\pi_0$}{initial guess}
        }
        {
            \inputitem{$\pi^*$}{optimal policy}
        }
        {
            \Repeat{convergence of $\pi$}
            {
                \ForEach{$s \in \mathcal{S}$}
                {
                    \ForEach{$a \in \mathcal{A}(s)$}
                    {
                        \tcp{policy evaluation}
                        collect sufficiently many episodes starting from $(s,a)$, following policy $\pi_k$\;
                        $q_{\pi_k}(s,a) \leftarrow$ the average return of all episodes starting from $(s,a)$\;
                    }
                }
                \tcp{policy improvement}
                $\pi_{k+1}(a|s) = \begin{cases}
                    1 &\text{if $a=\arg\max_{a'}{q_{\pi_k}(s,a')}$} \\
                    0 &\text{otherwise}
                \end{cases}$\;
            }
            \Return{$\pi$}
        }
    \subsection{Convergence}
    \subsection{Pros and Cons}
        Cons:
        \begin{itemize}
            \item Very inefficient sample usage. Due to this property, MC-Basic is not used in practical.
                To deal with this, we use MC with exploring start, which is presented in next section.
        \end{itemize}